\section{Arquitectura del sistema desarrollado}
\label{sec:arquitectura}

\subsection{Visi\'on general}

El sistema es una red inal\'ambrica de nodos de adquisici\'on s\'ismica coordinada por un
nodo maestro. La cadena completa, desde el suelo hasta el perfil $V_S$, es:

\begin{center}\footnotesize
\fbox{\parbox{0.95\textwidth}{\centering
Ge\'ofono SM-24 $\rightarrow$ cadena anal\'ogica (PSoC~5LP) $\rightarrow$ ADC $\Delta\Sigma$
$\rightarrow$ filtro FIR por hardware + DMA $\rightarrow$ RAM / SD \\[0.2em]
$\downarrow$ UART \\[0.2em]
ESP32 esclavo $\rightarrow$ ESP-NOW $\rightarrow$ ESP32 maestro (AP WiFi + SPA web + WebSocket) \\[0.2em]
$\downarrow$ \\[0.2em]
Servidor Python (FastAPI) $\rightarrow$ revisi\'on de campo $\rightarrow$ MASW $\rightarrow$ inversi\'on $\rightarrow$ $V_S(z)$
}}
\end{center}

Existen dos tipos de nodo, que comparten firmware: el \textbf{nodo GEO} (ge\'ofono) y el
\textbf{nodo HAMMER} (martillo instrumentado, que provee el instante de disparo y la
referencia de la fuente). El nodo maestro no adquiere: coordina, sincroniza y expone la
interfaz de operador.

\subsection{Nodo ge\'ofono: cadena de acondicionamiento anal\'ogico}

La cadena anal\'ogica, derivada de la topolog\'ia propuesta por Ma \emph{et
al.}~\cite{Ma2023}, se implementa \'integramente con los bloques anal\'ogicos configurables
internos del PSoC~5LP y componentes externos discretos. Consta de cuatro etapas en
cascada:

\begin{enumerate}[itemsep=0.15em]
    \item \textbf{PGA} --- amplificador de instrumentaci\'on de ganancia programable, en
    topolog\'ia de dos amplificadores operacionales, con ganancias seleccionables
    $\times1$, $\times4$, $\times16$, $\times32$ y $\times50$.
    \item \textbf{BP} --- filtro pasa-banda, que define la banda \'util y rechaza tanto la
    deriva de continua como el contenido de alta frecuencia fuera de inter\'es.
    \item \textbf{SUM} --- etapa sumadora inversora, que combina se\~nal y referencia.
    \item \textbf{LP} --- filtro pasa-bajos de segundo orden, que act\'ua como
    \emph{anti-aliasing} y alimenta directamente el ADC.
\end{enumerate}

Cada etapa dispone de un \textbf{VDAC dedicado} que fija su punto de operaci\'on. Cada
nodo de referencia se desacopla mediante una red \SI{1}{\micro\farad}\,$\|$\,\SI{100}{\nano\farad}
que, junto con la resistencia de salida interna de \SI{16}{\kilo\ohm} del VDAC, forma un
pasa-bajos de un polo en aproximadamente \SI{9}{\hertz}, atenuando la densidad de ruido
de \SI{750}{\nano\volt\per\sqrt{\hertz}} del VDAC por encima de esa frecuencia.

\subsection{\emph{Front-end} anal\'ogico asistido digitalmente (DA-AFE)}
\label{sec:arq-daafe}

El problema que motiva esta parte del dise\~no es la acumulaci\'on de \emph{offset}: en una
cadena de cuatro etapas con ganancia elevada, los \emph{offsets} intr\'insecos de cada
amplificador y de cada referencia se propagan y amplifican, consumiendo rango din\'amico
y, en el peor caso, saturando etapas intermedias antes de que llegue la se\~nal s\'ismica.

La soluci\'on adoptada es un \textbf{Digitally Assisted Analog Front-End}: durante una
fase de calibraci\'on previa a la adquisici\'on, los recursos digitales embebidos
---multiplexor anal\'ogico, ADC $\Delta\Sigma$, filtro digital programable, VDACs y
procesador--- se reconfiguran como un sistema de realimentaci\'on de se\~nal mixta que
recentra cada etapa. Terminada la calibraci\'on, esos recursos vuelven a su funci\'on de
adquisici\'on, los c\'odigos de los VDAC quedan fijos y el camino de se\~nal opera de forma
convencional, sin ning\'un lazo digital activo.

\subsubsection{Modelo}

Bajo las hip\'otesis de \emph{offsets} cuasi-est\'aticos durante la calibraci\'on, carga
despreciable entre etapas y superposici\'on lineal de las contribuciones de continua, el
\emph{offset} residual a la salida de la etapa $i$ se expresa como

\begin{equation}
y_i^{\mathrm{DC}}=\sum_{j<i}A_{ij}\,y_j^{\mathrm{DC}}+b_i+K_i\,u_i,
\label{eq:stage}
\end{equation}

donde $A_{ij}$ es la ganancia de continua de la etapa $j$ a la etapa $i$, $b_i$ el
\emph{offset} intr\'inseco de la etapa $i$, y $K_i$ la ganancia que relaciona la tensi\'on
inyectada por el VDAC $u_i$ con la salida de esa etapa. La estructura triangular de
\eqref{eq:stage} justifica el orden de calibraci\'on: de PGA a LP, porque cada etapa
hereda el residuo propagado por las anteriores.

\subsubsection{Ley de control}

El error se mapea primero al dominio del DAC de 8 bits:

\begin{equation}
\varepsilon_i[n]=\operatorname{round}\!\left(
\frac{e_i[n]\,V_{\mathrm{ADC,span}}\,N_{\mathrm{DAC}}}
{N_{\mathrm{ADC}}\,V_{\mathrm{DAC,span}}}\right),
\label{eq:scale}
\end{equation}

con $N_{\mathrm{ADC}}=2^{18}$, $N_{\mathrm{DAC}}=2^{8}$,
$V_{\mathrm{ADC,span}}=\SI{5000}{\milli\volt}$ y
$V_{\mathrm{DAC,span}}=\SI{4096}{\milli\volt}$.

Para evitar oscilaci\'on entre c\'odigos adyacentes cuando el punto de operaci\'on ideal no
es realizable, se adopta una banda muerta con piso conservador
$\Delta_{i,\min}^{\mathrm{impl}}=\max(1,\lceil|K_i|\rceil)$, dentro de la cual se
suprime la integraci\'on. El controlador es un \textbf{PI en forma posicional}:

\begin{equation}
c_i^{*}[n]=c_{i,\mathrm{base}}-s_i\left[
\frac{k_{\mathrm{P},i}}{|K_i|}\bar{\varepsilon}_i[n]
+\frac{k_{\mathrm{I},i}}{|K_i|}\sum_{q=0}^{n-1}\bar{\varepsilon}_i[q]\right],
\label{eq:pi}
\end{equation}

con $c_{i,\mathrm{base}}$ el c\'odigo base de prealimentaci\'on y
$s_i=\operatorname{sgn}(K_i)$. El comando se satura al rango 0--255 del DAC, con
\textbf{anti-windup direccional}: la integraci\'on se congela \'unicamente cuando el
residuo empujar\'ia el comando m\'as adentro de la saturaci\'on. Las ganancias $|K_i|$ se
obtuvieron por an\'alisis simb\'olico de continua de cada etapa en MATLAB a partir de los
valores nominales de componentes.

\begin{table}[h]
\centering\footnotesize
\caption{Par\'ametros de calibraci\'on empleados en todas las mediciones reportadas. Las
bandas muertas est\'an en c\'odigos equivalentes de DAC; S/L/R son muestras de
establecimiento/enganche/refinamiento a \SI{2604}{\hertz}, la tasa vigente durante esa campa\~na de medici\'on.}
\label{tab:cal_params}
\begin{tabular}{lccccccc}
\toprule
Etapa & $|K_i|$ & $\Delta_{i,\min}^{\mathrm{impl}}$ & $\Delta_i$ & $k_{\mathrm{P},i}$ & $k_{\mathrm{I},i}$ & S/L/R & \emph{Timeout} (s)\\
\midrule
PGA & 1    & 1 & 18 & $10^{-3}$ & $3\times10^{-4}$ & 512/512/1024   & 11,5\\
BP  & 1    & 1 & 3  & $10^{-4}$ & $5\times10^{-4}$ & 512/1024/1024  & 17,3\\
SUM & 7,89 & 8 & 10 & $10^{-4}$ & $5\times10^{-4}$ & 512/512/1024   & 11,5\\
LP  & 6    & 6 & 8  & $10^{-4}$ & $5\times10^{-4}$ & 1024/1024/2048 & 17,3\\
\bottomrule
\end{tabular}
\end{table}

\subsubsection{Flujo y persistencia}

Al arranque, tras un cambio de ganancia o por comando, el firmware carga los c\'odigos de
DAC almacenados y verifica si la calibraci\'on previa sigue siendo v\'alida; si lo es,
retorna inmediatamente a adquisici\'on. En caso contrario calibra las cuatro etapas
secuencialmente, con reinicio del filtro y establecimiento entre etapas. Un
\emph{watchdog} y l\'imites de conteo de muestras acotan el proceso.

Durante la calibraci\'on el filtro digital programable opera como estimador de continua
FIR de 128 \emph{taps}; durante la adquisici\'on se reconfigura con coeficientes adaptados
al ancho de banda del ge\'ofono. Un registro en EEPROM \textbf{indexado por c\'odigo de
ganancia} almacena la m\'ascara de etapas v\'alidas, los cuatro c\'odigos de DAC, la versi\'on
y un CRC; las calibraciones fallidas conservan sus mejores c\'odigos acotados sin
sobrescribir un registro v\'alido.

\subsection{Adquisici\'on digital en el PSoC~5LP}

\subsubsection{Cadena ADC--DMA--Filtro}

La adquisici\'on evolucion\'o desde una ISR cl\'asica por muestra hacia una cadena
gobernada por DMA en la que el procesador ARM Cortex-M3 no interviene en el camino de
datos. El bloque de filtro digital programable (DFB) aplica el FIR de adquisici\'on, y un
conjunto de canales DMA transporta las muestras del ADC al filtro y del filtro a memoria.

La \textbf{frecuencia de muestreo nativa del sistema es \SI{2929}{\hertz}}. Este valor no
es arbitrario: es la tasa real que resulta de la configuraci\'on de reloj del PSoC operando
el ADC $\Delta\Sigma$ a su \textbf{resoluci\'on m\'axima de 18~bits}. Es decir, la
arquitectura prioriza resoluci\'on sobre tasa de muestreo, y la frecuencia efectiva es
consecuencia de esa elecci\'on y no un par\'ametro fijado libremente.

A partir de esa tasa nativa se submuestrea por decimaci\'on con promediado, con factor $N$
configurable de 1 a 100:

\begin{equation}
F_s=\frac{\SI{2929}{\hertz}}{N}.
\end{equation}

El promediado de $N$ muestras consecutivas aporta, adem\'as de la reducci\'on de tasa, una
mejora de la relaci\'on se\~nal--ruido para el ruido incorrelado. El par\'ametro $N$ permite
intercambiar ancho de banda por duraci\'on de captura y por relaci\'on se\~nal--ruido, y se
configura desde la interfaz web. El ADC ofrece adem\'as cuatro configuraciones de rango de
entrada seleccionables, todas validadas en hardware.

\nota{\textbf{Punto a unificar}: los par\'ametros de calibraci\'on de la
Tabla~\ref{tab:cal_params} y los tiempos reportados en la Secci\'on~\ref{sec:resultados}
est\'an referidos a \SI{2604}{\hertz}, que es la tasa con la que se realiz\'o la campa\~na de
medici\'on del art\'iculo URUCON. La configuraci\'on vigente es \SI{2929}{\hertz}. Hay que
decidir si se re-expresan esos tiempos a la tasa actual (cambiar\'ian los valores ya
publicados) o si se aclara expl\'icitamente que corresponden a la configuraci\'on de aquel
experimento. Mi sugerencia es lo segundo: no tocar n\'umeros ya sometidos a revisi\'on
externa, y declarar la tasa junto a ellos.}

\subsubsection{\emph{superMaquina}: la m\'aquina de estados en hardware}

La l\'ogica de captura se migr\'o de C puro (ISRs sobre un registro de selecci\'on) a un
componente descrito en \textbf{Verilog} e instanciado en los bloques UDB del PSoC,
denominado \emph{superMaquina}. Este bloque decide en hardware qu\'e solicitud de DMA se
deja pasar, cu\'ando arranca y termina una captura, y qu\'e interrupciones llegan al
procesador.

La motivaci\'on es el determinismo: en un sistema multicanal el \emph{jitter} temporal se
traduce directamente en error de velocidad de fase. Al sacar la decisi\'on de captura del
dominio del software se elimina la dependencia de la latencia de interrupci\'on. La
pol\'itica de eventos asociada establece que los \emph{Terminal Count} de los temporizadores
fijos no generan interrupci\'on ---quedan retenidos en su registro de estado y los consume
el \emph{polling} del bucle de servicio---, y que durante el muestreo el nodo no atiende
RX/TX ni parpadeo de LED.

\subsubsection{Jerarqu\'ia de memoria y almacenamiento en SD}

Para soportar capturas largas se implement\'o una jerarqu\'ia de tres niveles:

\begin{enumerate}[itemsep=0.15em]
    \item \textbf{RAM del ESP32}: arena paginada de 512 lotes, el camino r\'apido para
    capturas cortas.
    \item \textbf{Encadenado en RAM del PSoC}: al superar los 512 lotes crudos, la
    \emph{superMaquina} se re-arma sin reiniciar, emitiendo un \'unico evento de fin al
    terminar.
    \item \textbf{Tarjeta SD en el PSoC}: sistema de archivos FatFs sobre el SPI maestro
    del TopDesign, con hasta 60\,000 lotes ($\approx\SI{11.5}{\minute}$ a
    \SI{2929}{\hertz}). Comandos UART dedicados para estado, autotest, captura a SD y
    lectura de lote.
\end{enumerate}

Es relevante destacar que la SD reside \textbf{en el PSoC y no en el ESP32}, lo que
mantiene el almacenamiento del lado determinista de la cadena.

\subsection{Red inal\'ambrica y sincronizaci\'on}

\subsubsection{Topolog\'ia}

Cada nodo de adquisici\'on integra un ESP32 conectado al PSoC por UART m\'as una l\'inea
dedicada de sincronismo. Los esclavos se comunican con el maestro por \textbf{ESP-NOW};
el maestro levanta simult\'aneamente un punto de acceso WiFi (\texttt{GeoNetwork},
\texttt{192.168.4.1}) operando en modo \texttt{WIFI\_AP\_STA}.

\subsubsection{Protocolo de disparo sincronizado}

La secuencia de disparo es el punto cr\'itico del sistema, ya que de ella depende la
alineaci\'on temporal entre canales:

\begin{enumerate}[itemsep=0.15em]
    \item \textbf{PRESTART}: el maestro difunde la configuraci\'on de captura; cada
    esclavo la transfiere al PSoC y lo deja armado.
    \item \textbf{HOT\_WAIT}: el maestro consulta a cada nodo si su PSoC confirm\'o el
    armado.
    \item \textbf{START\_NOW}: con todos los nodos armados, el maestro emite el disparo;
    la captura arranca por \textbf{flanco de una l\'inea f\'isica de sincronismo}, no por
    comando de software.
    \item \textbf{Dump}: finalizada la captura, los datos se transfieren mediante un
    esquema \emph{stop-and-wait} de solicitud de lotes, con pausa y reanudaci\'on seg\'un
    la presencia del cliente WebSocket.
\end{enumerate}

El arranque por flanco f\'isico es la decisi\'on de dise\~no que hace viable la
sincronizaci\'on multicanal: elimina de la ruta cr\'itica la latencia variable del
\emph{stack} de radio y del software.

\subsection{Nodo martillo (fuente s\'ismica)}

El nodo HAMMER acondiciona un martillo instrumentado con aceler\'ometro
PCB~Piezotronics~086D20 y provee el instante de disparo y la referencia de la se\~nal de
la fuente. Comparte el \emph{mismo archivo de firmware} que el nodo GEO: el c\'odigo
detecta autom\'aticamente en tiempo de compilaci\'on qu\'e tipo de hardware tiene debajo
seg\'un los componentes presentes en el TopDesign, y configura la cadena anal\'ogica y los
objetivos de calibraci\'on correspondientes. Esta decisi\'on elimina la bifurcaci\'on del
firmware embebido y garantiza que el protocolo y la l\'ogica de captura sean id\'enticos en
ambos tipos de nodo. La clase de hardware se propaga desde el PSoC hasta la interfaz web.

\subsection{Interfaz de operador en campo}

El maestro sirve una aplicaci\'on web de p\'agina \'unica desde su propio sistema de archivos
LittleFS sobre su AP WiFi. Esto elimina la necesidad de llevar una computadora al campo:
el operador se conecta con el tel\'efono. La interfaz permite configurar decimaci\'on,
duraci\'on por tipo de nodo, l\'imites de RAM y SD, y rango de ADC; muestra el estado de
cada nodo con indicador de intensidad de se\~nal (RSSI) del enlace ESP-NOW; refleja las
capturas en vivo por WebSocket; y permite exportar la sesi\'on como archivo comprimido
armado \'integramente en el navegador.

\subsection{\emph{Pipeline} de procesamiento MASW}
\label{sec:arq-pipeline}

El procesamiento posterior se implement\'o en Python, sobre un servidor FastAPI que
reemplaz\'o a una primera versi\'on artesanal y a una aplicaci\'on de escritorio PyQt previa.
Sus componentes son:

\begin{itemize}[itemsep=0.2em]
    \item \textbf{Revisi\'on de campo}: inspecci\'on de capturas con ordenamiento por
    amplitud pico a pico, banco de filtros configurable, alineaci\'on entre tandas
    (\emph{enfase}) por carpeta, correcci\'on de polaridad por punto de medici\'on
    (autom\'atica y manual), mezcla de capturas tomadas a distintas frecuencias de
    muestreo mediante remuestreo racional, deduplicaci\'on por captura, promediado por
    carpeta y rechazo expl\'icito de carpetas defectuosas.
    \item \textbf{An\'alisis de dispersi\'on}: construcci\'on de la imagen MASW en el dominio
    frecuencia--velocidad de fase, visualizaci\'on tipo \emph{waterfall} y $f$--$k$, y
    definici\'on de \textbf{regiones-pol\'igono} que permiten extraer varias curvas
    simult\'aneas, una por modo.
    \item \textbf{Inversi\'on}: inversi\'on multimodal usando \texttt{evodcinv} sobre el
    motor de modelado directo \texttt{disba}, con la curva editable antes de invertir y
    persistencia de los resultados intermedios.
    \item \textbf{Infraestructura}: cola de trabajos de an\'alisis persistente, borrado
    por cuarentena reversible en lugar de eliminaci\'on f\'isica, exportaciones
    reproducibles y un \emph{gate} automatizado de 58 verificaciones como condici\'on de
    avance.
\end{itemize}

\subsection{Organizaci\'on del repositorio}

El proyecto se reorganiz\'o desde un repositorio monol\'itico hacia un superproyecto con
subm\'odulos independientes, agrupados \textbf{por prop\'osito y no por lenguaje}:

\begin{center}\footnotesize
\begin{tabular}{@{}ll@{}}
\toprule
\texttt{src/firmware/psoc} & Firmware y dise\~no anal\'ogico del PSoC~5LP \\
\texttt{src/firmware/esp32} & Firmware de los nodos maestro y esclavo \\
\texttt{src/interfaces/python} & Servidor FastAPI, revisi\'on de campo, exportaci\'on \\
\texttt{src/interfaces/matlab} & Interfaces de banco (osciloscopio, ESP) \\
\texttt{src/calculos\_modelados/matlab} & Modelo SM-24, compensadores, an\'alisis de circuito \\
\texttt{src/calculos\_modelados/python} & MASW, martillo de leva \\
\texttt{data} & Mediciones crudas y procesadas (Git LFS, dedup.\ por SHA-256) \\
\texttt{docs} & Documentaci\'on, bit\'acora, comunicaciones \\
\texttt{PCBs} & Dise\~nos electr\'onicos (KiCad, JitX) \\
\bottomrule
\end{tabular}
\end{center}

Un \emph{commit} del superproyecto fija un SHA concreto de cada subm\'odulo, constituyendo
una l\'inea base reproducible. Los datos de medici\'on se versionan con Git~LFS con
deduplicaci\'on por SHA-256, conservando el repositorio los punteros, el cat\'alogo y la
estructura.

\nota{En la versi\'on de 15 p\'aginas esta secci\'on es la principal candidata a recorte
agresivo: el Art.~19 pide un documento de \emph{Estado del Arte}, no una memoria de
implementaci\'on. Sugerencia: reducir a ~2,5 p\'aginas conservando \S4.1 (diagrama),
\S4.2, \S4.3 (que es el aporte metodol\'ogico y ya tiene revisi\'on externa), \S4.5.2
(sincronizaci\'on, que es el problema t\'ecnico distintivo) y \S4.8. Las secciones
\S4.4.3, \S4.6, \S4.7 y \S4.9 pueden resumirse en un p\'arrafo cada una o pasar a la
Defensa Final.}
